\chapter{Zadatak 1 --- priprema za obranu}
\label{ch:z1-obrana}

\section{Brojevi koje moraš znati napamet}
\label{sec:z1o-brojevi}

Ovo je jedina tablica u udžbeniku koju treba naučiti doslovno. Sve ostalo se
izvodi.

\begin{longtable}{@{}p{0.30\linewidth} p{0.30\linewidth} p{0.34\linewidth}@{}}
\toprule
\textbf{Veličina} & \textbf{Egzaktno} & \textbf{Decimalno} \\
\midrule
\endhead
$\lambda(Q)$ & --- & $1{,}955969$ \ i \ $4{,}044031$ \\
$\det Q$ & $791/100$ & $7{,}91$ \\
$x^{0}$ (neogr.\ min.) & $(-710/791,\; 540/791)$ & $(-0{,}897598,\; 0{,}682680)$ \\
$g_1(x^{0})$ & $8080/791$ & $10{,}214918$ \\
$\xopt$ & $(190/33,\; 140/33)$ & $(5{,}757576,\; 4{,}242424)$ \\
$\muopt$ & $(404/33,\; 0)$ & $(12{,}242424,\; 0)$ \\
$\popt = \dopt$ & $2000/33$ & $60{,}606061$ \\
$g_2(\xopt)$ (rezerva) & $-49/33$ & $-1{,}484848$ \\
$\grad f(\xopt)$ & $(404/33,\; 404/33)$ & $(12{,}242424,\; 12{,}242424)$ \\
Vrh dopustivog skupa & $(13/2,\; 7/2)$ & $(6{,}5,\; 3{,}5)$ \\
Množitelji u vrhu & $(23/2,\; -49/20)$ & $(11{,}50,\; -2{,}45)$ \\
Slaterova točka & $(7,\; 6)$ & $g_1 = -3$, $g_2 = -2$ \\
$\partial\popt/\partial v$ & $(-404/33,\; 0)$ & $(-12{,}242424,\; 0)$ \\
\bottomrule
\end{longtable}

\begin{izgovor}[Trik za pamćenje]
Sve što je važno ima nazivnik $33$. $x_1^{\star} = 190/33$, $x_2^{\star} = 140/33$,
$\mu_1^{\star} = 404/33$, $\popt = 2000/33$, rezerva $= -49/33$. Zapamti pet
brojnika --- $190$, $140$, $404$, $2000$, $-49$ --- i nazivnik $33$, i imaš cijeli
zadatak. Provjera: $190 + 140 = 330 = 10\cdot33$, što je upravo aktivno
ograničenje $x_1 + x_2 = 10$.
\end{izgovor}

\section{Zadatak u pet rečenica}
\label{sec:z1o-sazetak}

\begin{izgovor}[Ako imaš samo minutu]
,,Zadatak 1 je kvadratni program s dva ograničenja nejednakosti. Problem je
konveksan --- Hesseova matrica je pozitivno definitna, ograničenja su afina --- i
Slaterov uvjet vrijedi, pa su KKT uvjeti nužni i dovoljni.

Riješio sam KKT sustav nabrajanjem sve četiri kombinacije aktivnosti. Tri
otpadaju: jedna zbog nedopustivosti, dvije zbog negativnog množitelja. Preostaje
slučaj u kojem je aktivno samo prvo ograničenje, s rješenjem sto devedeset kroz
trideset tri i sto četrdeset kroz trideset tri.

\texttt{quadprog} daje isto do na jedanaestu decimalu, i vraća iste množitelje.

Dualni problem sam izveo analitički --- opet kvadratni program, u varijabli mi --- i
riješio numerički; dualni procjep je numerička nula, kako i mora biti uz jaku
dualnost.

Osjetljivost slijedi izravno iz množitelja: samo prvo ograničenje ima cijenu,
dvanaest zarez dva četiri po jedinici; drugo je neaktivno i nema utjecaja dok
perturbacija ne premaši rezervu od jedan zarez četiri osam.``
\end{izgovor}

\section{Pitanja koja se očekuju}
\label{sec:z1o-pitanja}

\subsection{O konveksnosti}

\begin{pitanje}{Kako biste dokazali konveksnost da funkcija cilja nije kvadratna?}
\textbf{Odgovor.} Kriterij drugog reda vrijedi i tada, samo Hessian više nije
konstantan, pa treba pokazati $\grad^2 f(x) \psd 0$ \emph{za sve $x$ u domeni}.
Ako to nije izvedivo, ostaje definicija preko tetive,
$f(\alpha x_1 + (1-\alpha)x_2) \leq \alpha f(x_1) + (1-\alpha)f(x_2)$, ili
sastavljanje funkcije od poznatih konveksnih gradivnih blokova pravilima koja ih
čuvaju --- zbroj, maksimum, kompozicija s afinim preslikavanjem.
\end{pitanje}

\begin{pitanje}{Što ako je $Q$ samo pozitivno semidefinitna, a ne definitna?}
\textbf{Odgovor.} Problem je i dalje konveksan, ali funkcija cilja nije strogo
konveksna, pa rješenje ne mora biti jedinstveno --- može postojati cijeli
optimalni segment. Osim toga, $Q^{-1}$ tada ne postoji, pa izvod dualne funkcije
iz podzadatka (e) u ovom obliku propada i treba ga voditi preko pseudoinverza,
uz uvjet dopustivosti na $c + G\trans\mu$.
\end{pitanje}

\begin{pitanje}{Zašto je bilo važno provjeriti Slaterov uvjet?}
\textbf{Odgovor.} Jer bez njega KKT uvjeti ostaju samo nužni. Iz nužnosti slijedi
da optimum \emph{ako postoji} zadovoljava KKT; ne slijedi da točka koja
zadovoljava KKT jest optimum. Slater daje jaku dualnost, a odatle i dovoljnost ---
teorem s Predavanja 4, slajd 17/26.
\end{pitanje}

\subsection{O KKT uvjetima}

\begin{pitanje}{Objasnite fizikalno uvjet komplementarnosti.}
\textbf{Odgovor.} Ograničenje koje ne dodiruje rješenje ne smije na njega ni
djelovati. Ako je $g_j(\xopt) < 0$, ograničenje je ,,daleko`` i njegova sila,
$\mu_j$, mora biti nula. Obrnuto, ako je sila različita od nule, ograničenje mora
biti aktivno --- rješenje leži na njemu. Umnožak $\mu_j g_j = 0$ zapisuje upravo
tu ekskluzivnost.
\end{pitanje}

\begin{pitanje}{Što bi značilo da je gradijent aktivnog ograničenja nul-vektor?}
\textbf{Odgovor.} Da uvjet regularnosti pada. Gradijenti aktivnih ograničenja
tada nisu linearno nezavisni, KKT teorem se ne smije primijeniti i množitelji ne
moraju postojati. Kod nas je $\grad g_1 = (-1,-1) \neq 0$, i to je jedini aktivan
gradijent, pa je linearna nezavisnost trivijalno zadovoljena.
\end{pitanje}

\begin{pitanje}{Zašto se u Slučaju 2 množitelj $\mu_1$ mogao eliminirati?}
\textbf{Odgovor.} Jer je u obje skalarne jednadžbe stacionarnosti ušao s istim
koeficijentom, $-1$, budući da su obje komponente $\grad g_1$ jednake $-1$.
Izjednačavanjem izraza za $\mu_1$ iz dviju jednadžbi množitelj ispada, i ostaje
jedna linearna veza između $x_1$ i $x_2$, koja se s uvjetom aktivnosti rješava
kao sustav $2\times2$.
\end{pitanje}

\subsection{O dualnosti}

\begin{pitanje}{Što bi značio pozitivan dualni procjep?}
\textbf{Odgovor.} Da jaka dualnost ne vrijedi --- dakle da problem nije konveksan
ili da nije zadovoljen uvjet kvalifikacije ograničenja. Kod nas je oboje
zadovoljeno, pa je procjep nula; izračunatih $2{,}48\cdot10^{-10}$ je tolerancija
rješavača.
\end{pitanje}

\begin{pitanje}{Koliko varijabli ima dualni problem?}
\textbf{Odgovor.} Onoliko koliko primar ima ograničenja --- ovdje dvije. Primar
ima dvije varijable i dva ograničenja, pa su iste veličine. Kod problema s mnogo
varijabli i malo ograničenja dual može biti bitno manji, i to je jedan od
praktičnih razloga zašto se uopće koristi.
\end{pitanje}

\subsection{O osjetljivosti}

\begin{pitanje}{Zašto u formuli osjetljivosti stoji minus?}
\textbf{Odgovor.} Zato što $v_i > 0$ popušta ograničenje i povećava dopustivi
skup, a minimum po većem skupu ne može porasti. Kako je $\mu_i^{\star} \geq 0$,
derivacija $-\mu_i^{\star}$ je nepozitivna --- popuštanje smanjuje ili ne mijenja
$\popt$.
\end{pitanje}

\begin{pitanje}{Vrijedi li rezultat za veliku perturbaciju?}
\textbf{Odgovor.} Kao derivacija --- ne, samo lokalno, dok se aktivni skup ne
promijeni. Ali globalna nejednakost
$\popt(v) \geq \popt(0) - \muopt{}\trans v$ vrijedi za svaki $v$, jer slijedi iz
slabe dualnosti primijenjene na perturbirani problem.
\end{pitanje}

\section{Popis zamki}
\label{sec:z1o-zamke}

\begin{longtable}{@{}p{0.05\linewidth} p{0.42\linewidth} p{0.47\linewidth}@{}}
\toprule
\textbf{\#} & \textbf{Zamka} & \textbf{Kako je izbjeći} \\
\midrule
\endhead
1 & Faktor $\tfrac12$ ispred kvadratne forme & Uvijek raspiši
$\tfrac12 x\trans Qx$ natrag i usporedi sa zadanom funkcijom (\S\ref{sec:kvadratna-forma}). \\[5pt]
2 & Dijeljenje mješovitog člana na pola & $Q_{12} = -0{,}3$, ne $-0{,}15$;
polovina je već u faktoru ispred. \\[5pt]
3 & Predznak pri prebacivanju $-10$ & $-x_1-x_2 \leq -10 \Leftrightarrow -x_1-x_2+10 \leq 0$
(\S\ref{sec:z1-ogranicenja}). \\[5pt]
4 & Odgovor samo o funkciji cilja & Zadatak pita je li \emph{problem} konveksan --- i
dopustivi skup mora u odgovor. \\[5pt]
5 & Zaključivanje o definitnosti iz predznaka elemenata & Samo svojstvene
vrijednosti, minori, ili $z\trans Qz$ (\S\ref{sec:relacije}). \\[5pt]
6 & Prešućena egzistencija minimuma & Skup je neograničen; treba spomenuti
koercivnost. \\[5pt]
7 & Rješavanje samo ,,pogođenog`` slučaja & Ispiši tablicu svih $2^m$ kombinacija,
kao profesor na Vježbama 3. \\[5pt]
8 & Negativan $\mu$ shvaćen kao greška u računu & To je rezultat --- dokaz da ta
kombinacija nije optimalna. \\[5pt]
9 & Zaboravljena provjera neaktivnog ograničenja & U Slučaju 2 treba pokazati i
$g_2(\xopt) < 0$. \\[5pt]
10 & Vrh dopustivog skupa uzet kao optimum & Ondje je $\mu_2 = -2{,}45 < 0$;
vrh otpada. \\[5pt]
11 & Elipsa nacrtana tako da siječe rub & U optimumu mora \emph{dodirivati}
(\S\ref{sec:z1c-crtanje}). \\[5pt]
12 & Krivi smjer $\grad f(\xopt)$ & Pokazuje \emph{u} dopustivi skup; van pokazuje
$-\grad f$. \\[5pt]
13 & Puna matrica bez $\tfrac12$ predana \texttt{quadprog}-u & Rješavač
podrazumijeva $\tfrac12$; predaje se $Q$. \\[5pt]
14 & Izostavljen konstantni član kod dualnog problema & Bez
$-\tfrac12 c\trans Q^{-1}c$ procjep ispadne $1{,}92$ umjesto nule. \\[5pt]
15 & ,,$\mu_2 = 0$ znači da ograničenje ne treba`` & Znači samo da nije aktivno;
rezerva je $1{,}4848$. \\[5pt]
16 & Izgubljen minus u osjetljivosti & Popuštanje smanjuje $\popt$; derivacija je
nepozitivna. \\
\bottomrule
\end{longtable}

\section{Redoslijed izlaganja pred pločom}
\label{sec:z1o-redoslijed}

\begin{ploca}[Kako izložiti cijeli Zadatak 1]
\begin{enumerate}[leftmargin=*]
  \item \textbf{Zapiši problem} u matričnom obliku, gore lijevo na ploči, i ostavi
        ga ondje --- na njega ćeš se pozivati cijelo vrijeme.
  \item \textbf{Konveksnost} u tri reda: $\grad^2 f = Q$, svojstvene vrijednosti
        $1{,}956$ i $4{,}044$, ograničenja afina. Dopiši Slaterovu točku $(7,6)$.
  \item \textbf{Tablica četiriju slučajeva} u desnom gornjem kutu. Prekriži tri,
        zaokruži jedan.
  \item \textbf{Slučaj 2 do kraja}: tri jednadžbe, eliminacija $\mu_1$,
        $6{,}6x_1 = 38$, rezultat. Uokviri $\xopt$ i $\muopt$.
  \item \textbf{Slika} u donjoj polovici ploče, po postupku iz
        \S\ref{sec:z1c-crtanje}. Ostavi joj dovoljno mjesta --- to je dio na kojem
        se najviše dobiva.
  \item \textbf{Dual} u dva reda: $\tilde{x} = -Q^{-1}(c + G\trans\mu)$, pa
        $\dual(\mu)$. Naglasi da je opet QP i da je $\dopt = \popt$.
  \item \textbf{Osjetljivost} u jednom redu: $\partial\popt/\partial v = -\muopt$,
        pa brojevi i rezerva.
\end{enumerate}
\end{ploca}

\begin{zamka}[Zamka: brisanje ploče]
Ne briši matrični zapis problema ni tablicu slučajeva. Ispitivač se vraća na njih
pitanjima tipa ,,a što bi bilo da je ovdje umjesto tri bilo pet``. Ako je ploča
prazna, počinješ ispočetka i gubiš vrijeme.
\end{zamka}

\section{Veze prema ostalim zadacima}
\label{sec:z1o-veze}

\begin{predznanje}
Što se iz Zadatka 1 nosi dalje:
\begin{itemize}[leftmargin=*]
  \item \textbf{U Zadatak 2}: isti problem, ali s nesigurnim koeficijentima
        $a_{ij}$. Robusna formulacija svodi se na konačno mnogo ograničenja, pa
        se opet rješava KKT uvjetima --- ista tablica aktivnosti, više redaka.
  \item \textbf{U Zadatak 3}: ista logika ,,funkcija cilja plus afina
        ograničenja``, samo je cilj $\ell_1$-norma, pa se pomoćnim varijablama
        svodi na linearni program.
  \item \textbf{U Zadatak 4}: pojam pozitivne definitnosti prelazi s Hesseove
        matrice na Ljapunovljevu matricu $P$; umjesto $Q \pd 0$ traži se
        $P \pd 0$ i $M \nd 0$, a to su linearne matrične nejednakosti.
\end{itemize}
\end{predznanje}
